\scanpage{16}
où les $\varphi_i$ sont des fonctions décroissantes \inserted{sommables} sur $]0,|K|]$ (\unsure{partie bornée dépendant de} $K$). Bien entendu, si $P$ est défini sur chaque $L^\infty(K)$ par une formule (1), $P$ est une fonction de Polya.

Supposons de plus que $P$ satisfasse :

$P_4)$ \uline{$P$ est croissante}

Alors on voit aussitôt que les $\varphi_i$ dans (1) sont positives (et réciproquement).

\deleted{En \unsure{prolongeant} \unsure{les} $\varphi_i$ par $0$ sur $[|K|,+\infty[$, et \unsure{considérant} \unsure{toutes les} $\varphi_i$ (pour tous les $K$) on obtient une formule unique
\[
(2)\qquad \boxed{P(f)=\mathop{\rm Sup}_i\ a_i+\int_0^\infty\varphi_i\varphi_f\,\unsure{\xi}}
\]
où les $\varphi_i$ sont une famille de fonctions positives décroissantes \inserted{sommables} sur $R^+$, \unsure{toutes nulles} \unsure{au voisinage de} $\infty$.
La formule (2) peut s'écrire
\[
(3)\qquad P(f)=P^*(\varphi_f)
\]
où $P^*$ est la fonction de Polya sur $R^+_*$ donnée par
\[
(4)\qquad P^*(\varphi)=\mathop{\rm Sup}_i\left(a_i+\int_0^\infty\varphi_i\varphi\right)
\]
}

\margin{\uline{Proposition 3} Soit $P$ une fonction de Polya sur $]0,\omega[$, (\deleted{$\omega$ fini} $0<\omega<+\infty$). Alors $P$ est de la forme (1), où les $\varphi_i$ sont des fonctions décroissantes $\in L^1(]0,\omega[)$. Pour que $P$ soit croissante, il faut et il suffit \unsure{qu'on puisse prendre} les $\varphi_i\geq0$.}

\margin{\deleted{Définition de $P(f)$ si $f\in\unsure{R^\infty}$} La prop. 3 \unsure{ci-dessus} \unsure{ne} $=$ \unsure{toujours la forme} \unsure{principale} des fonctions de Polya.}
\mspage{17}
